% ==============================================================================
% TEMPLATE: BAB DAFTAR PUSTAKA
% ==============================================================================
% File ini memuat struktur bab Daftar Pustaka menggunakan lingkungan \begin{enumerate}
% dengan format penomoran dan label \item \label{...} untuk disitasi via [\ref{...}].
% Ganti teks placeholder di bawah dengan referensi literatur riil.
% ==============================================================================

\chapter{Daftar Pustaka}
\label{chap:daftarpustaka}

\begin{enumerate}
    \item \label{modulPraktikum} [Nama Dosen Pengampu, Gelar], dan Tim Asisten. \textit{Modul Praktikum [Nama Mata Kuliah]: [Topik Praktikum Pertemuan X]}. Yogyakarta: Departemen Teknik Elektro dan Informatika, Sekolah Vokasi, Universitas Gadjah Mada, [Tahun].
    
    \item \label{bukuTeks} [Nama Penulis Buku, Inisial]. \textit{[Judul Buku Referensi: Subjudul]}, [Edisi] ed. [Kota Penerbit]: [Nama Penerbit], [Tahun].
    
    \item \label{jurnalIlmiah} [Penulis 1] and [Penulis 2], "[Judul Artikel Ilmiah / Riset Terkait]," \textit{[Nama Jurnal Ilmiah]}, vol. [Volume], no. [Nomor], pp. [Awal]--[Akhir], [Tahun]. [Online]. Tersedia: \url{https://doi.org/...}.
    
    \item \label{sumberDataset} [Penyedia / Penulis Dataset], \textit{[Nama Dataset / Kumpulan Data]}. Kaggle / Repositori Data, [Tahun]. [Online]. Tersedia: \url{https://...}.
    
    \item \label{sourceCode} [Nama Mahasiswa], \textit{[Kode Matkul]-[NIM] (Source Code Laporan Praktikum)}. Google Colaboratory / GitHub, [Tahun]. [Online]. Tersedia: \url{https://colab.research.google.com/drive/...}.
\end{enumerate}
